\chapter*{Resumen}
\addcontentsline{toc}{chapter}{Resumen}

El monitoreo sismo-volcánico automático requiere detectar eventos en
registro continuo y clasificarlos por tipo. Los sistemas en operación
resuelven ambas etapas de forma desacoplada: un detector fija los límites
del evento y un clasificador entrenado sobre segmentaciones hechas por
analistas asigna la clase. Esta tesis desarrolla un algoritmo de
identificación de eventos sismo-volcánicos que ajusta la etapa de
clasificación a la variabilidad del recorte que produce el detector, sobre
registros del Complejo Volcánico Nevados de Chillán.

El trabajo evalúa empíricamente tres detectores (STA/LTA, PhaseNet y
EQTransformer) en dos niveles, identificación del instante del evento y
calidad de la delimitación temporal; construye un procedimiento de recorte
adaptativo que se ajusta a la duración real de los eventos, cuyo rango
excede en dos órdenes de magnitud la ventana fija empleada habitualmente; y
valida el algoritmo integrado contrastando arquitecturas convolucionales
alternativas y cuantificando la degradación entre el recorte de referencia
y el recorte automático.

\vspace{1cm}
\noindent\textbf{Palabras clave:} sismología volcánica, detección de
eventos, clasificación automática, redes neuronales convolucionales,
Nevados de Chillán.